\answer{ex:20:1}{$t_{\text{нс}}=\tau_r\ln\frac{1-L}{1-H}=16.8$~год, $t_{\text{с}}=\tau_d\ln\frac HL=5.8$~год, період $22.5$~год; до $24$~год генератор підтягує добове гойдання порогів --- фазове автопідстроювання~\eqref{eq:pll}.}
\answer{ex:20:2}{(а)~$L=0.111$; за $L=0.092$ сон тривав би $9.6$~год. (б)~На $22\%$ ($1/0.82=1.22$), а не на $18\%$.}
\answer{ex:20:3}{$H=\frac{p-1}{p-1/q}=0.623$, $L=H/q=0.093$, де $p=e^{16/\tau_r}$, $q=e^{8/\tau_d}$. Після пробудження через $4$~год фаза зсувається назавжди (засинання на $7.2$~год раніше); із гойдальними порогами вона повертається за дві--три доби.}
\answer{ex:20:4}{(а)~$r_\pm=\frac12\bigl(1\pm\sqrt{1-4k/w}\bigr)=0.947,\ 0.053$; $a_{\text{в}}=3.51$, $a_{\text{н}}=0.49$. (б)~Робота $\approx0.33$~с, тиша $\approx0.59$~с, період $\approx0.93$~с, частка роботи $0.36$.}
\answer{ex:20:5}{$a_{\text{в}}/r_+<b<a_{\text{н}}/r_-$: $3.71<b<9.20$. \nt{ACET} зменшує $b$ нижче за $3.71$: перетин переходить на верхню гілку, і стійким стає стан роботи ($r\approx1$).}
\answer{ex:20:6}{$4.9$, $6.8$, $8.8$, $5.5$, $8.9$~год за $D=24$, $32$, $40$, $48$, $64$: тривалість задає добова фаза засинання, а не борг.}
\answer{ex:20:7}{Після B похибка на A в середньому $0.51$ (від $0.19$ до $1.08$ за $20$ наборами; у нульової відповіді --- $1$). Упереміш обидві похибки менші за $10^{-4}$.}
\answer{ex:20:8}{Відкрита задача. Рівномірне масштабування зберігає відношення ваг, але зберігає відповіді порогового нейрона лише тоді, коли поріг масштабується в ту саму кількість разів; перемішані повтори змінюють відношення. Незмінність великих контактів суперечить суто рівномірному масштабуванню.}
